% GENERATED FILE. Edit canonical lessons, then run npm run generate:books.
% canonical-source-hash: fnv1a64:376ca26b9be9407e
% canonical-lessons: ZH-C338-guojie, ZH-C338-baifang, ZH-C338-tuanyuan, ZH-C338-jiari, ZH-C338-fengsu

\chapter{Festival Visits}
\label{ch:zh-festival-visits}
\hlchaptermodality{338}

\begin{chapteropening}
\textbf{By the end of this chapter:} \emph{I can talk about spending a festival at home, paying visits, family reunions, days off and local customs.}

Complete the last lesson of chapter 338: I can talk about spending a festival at home, paying visits, family reunions, days off and local customs.
\end{chapteropening}

\section[\texorpdfstring{\zh{过节} (guòjié)}{过节 (guòjié)}]{\zh{过节} (guòjié) — to celebrate a festival}

\label{lesson:ZH-C338-guojie}



\begin{quote}
Before the new one: say the Mandarin for good luck, then the Mandarin for to make a wish.
\end{quote}

\subsection*{You'll want to know: \zh{过节}}
\textbf{\zh{过节}} — \emph{guòjié} — \textquotedblleft{}to celebrate a festival\textquotedblright{}.

\begin{quote}
\textbf{\zh{去年我们回老家过节}} \emph{qùnián wǒmen huí lǎojiā guòjié} — Last year we went back to our home town for the festival.
\end{quote}

\begin{cousinweb}
\textbf{\zh{过节}} (\emph{guòjié}) is built from \zh{过} 'to pass' and \zh{节} 'a festival', both inherited from Middle Chinese, \emph{kwaH} and \emph{tset} (Guangyun, in Baxter's transcription), which glosses the first 'to go beyond' and the second 'a joint of bamboo'. The sense 'a festival' is not given there.
\end{cousinweb}

\begin{grammarlens}[title={What you've built}]
One more word for telling what happened.
\end{grammarlens}

\subsection*{Your turn}
\begin{itemize}
\raggedright
  \item \emph{Say it:} \emph{guòjié}
  \item \emph{Say it:} \emph{guòjié}, once more
  \item \emph{Say it:} \emph{qùnián wǒmen huí lǎojiā guòjié} — Last year we went back to our home town for the festival.
  \item \emph{Recall:} say the Mandarin for good luck, then the Mandarin for to make a wish, then say \emph{guòjié} again
\end{itemize}

\begin{tcolorbox}[breakable,colback=teal!4,colframe=teal!35!black,title={Before you move on}]
What does \zh{过节} mean? (\textquotedblleft{}To celebrate a festival\textquotedblright{}.) Say it once more.
\end{tcolorbox}
\section[bàifǎng]{bàifǎng — to pay a visit to}

\label{lesson:ZH-C338-baifang}



\begin{quote}
Before the new one: say the Mandarin for to make a wish, then the Mandarin for to celebrate a festival.
\end{quote}

\subsection*{You'll want to know: bàifǎng}
\emph{bàifǎng} — \textquotedblleft{}to pay a visit to\textquotedblright{}.

\begin{quote}
\emph{guòjié de shíhou wǒmen qù bàifǎng le lǎoshī} — Over the festival we went to visit our teacher.
\end{quote}

Its characters come later; for now it is a word for the ear.

\begin{cousinweb}
\emph{bàifǎng} is built from \emph{bài} 'to bow' and \emph{fǎng} 'to call on', both inherited from Middle Chinese, \emph{peajH} and \emph{phjangH} (Guangyun, in Baxter's transcription), which glosses the second 'to consult'.
\end{cousinweb}

\begin{grammarlens}[title={What you've built}]
One more word for telling what happened.
\end{grammarlens}

\subsection*{Your turn}
\begin{itemize}
\raggedright
  \item \emph{Say it:} \emph{bàifǎng}
  \item \emph{Say it:} \emph{bàifǎng}, once more
  \item \emph{Say it:} \emph{guòjié de shíhou wǒmen qù bàifǎng le lǎoshī} — Over the festival we went to visit our teacher.
  \item \emph{Recall:} say the Mandarin for to make a wish, then the Mandarin for to celebrate a festival, then say \emph{bàifǎng} again
\end{itemize}

\begin{tcolorbox}[breakable,colback=teal!4,colframe=teal!35!black,title={Before you move on}]
What does bàifǎng mean? (\textquotedblleft{}To pay a visit to\textquotedblright{}.) Say it once more.
\end{tcolorbox}
\section[tuányuán]{tuányuán — to be reunited as a family}

\label{lesson:ZH-C338-tuanyuan}



\begin{quote}
Before the new one: say the Mandarin for to celebrate a festival, then the Mandarin for to pay a visit to someone.
\end{quote}

\subsection*{You'll want to know: tuányuán}
\emph{tuányuán} — \textquotedblleft{}to be reunited as a family\textquotedblright{}.

\begin{quote}
\emph{guònián de shíhou jiārén dōu huílai tuányuán} — At the new year the whole family comes back together.
\end{quote}

Its characters come later; for now it is a word for the ear.

\begin{cousinweb}
\emph{tuányuán} is built from \emph{tuán} and \emph{yuán}, both inherited from Middle Chinese, \emph{dwan} and \emph{hjwen} (Guangyun, in Baxter's transcription), which glosses the first with this very word and the second 'the round body of the heavens'.
\end{cousinweb}

\begin{grammarlens}[title={What you've built}]
One more word for telling what happened.
\end{grammarlens}

\subsection*{Your turn}
\begin{itemize}
\raggedright
  \item \emph{Say it:} \emph{tuányuán}
  \item \emph{Say it:} \emph{tuányuán}, once more
  \item \emph{Say it:} \emph{guònián de shíhou jiārén dōu huílai tuányuán} — At the new year the whole family comes back together.
  \item \emph{Recall:} say the Mandarin for to celebrate a festival, then the Mandarin for to pay a visit to someone, then say \emph{tuányuán} again
\end{itemize}

\begin{tcolorbox}[breakable,colback=teal!4,colframe=teal!35!black,title={Before you move on}]
What does tuányuán mean? (\textquotedblleft{}To be reunited as a family\textquotedblright{}.) Say it once more.
\end{tcolorbox}
\section[\texorpdfstring{\zh{假日} (jiàrì)}{假日 (jiàrì)}]{\zh{假日} (jiàrì) — holiday; day off}

\label{lesson:ZH-C338-jiari}



\begin{quote}
Before the new one: say the Mandarin for to pay a visit to someone, then the Mandarin for to be reunited as a family.
\end{quote}

\subsection*{You'll want to know: \zh{假日}}
\textbf{\zh{假日}} — \emph{jiàrì} — \textquotedblleft{}holiday; day off\textquotedblright{}.

\begin{quote}
\textbf{\zh{祝你假日快乐}} \emph{zhù nǐ jiàrì kuàilè} — Have a happy holiday!
\end{quote}

\begin{cousinweb}
\textbf{\zh{假日}} (\emph{jiàrì}) is built from \zh{假} 'leave' and \zh{日} 'a day', both inherited from Middle Chinese, \emph{kaeH} and \emph{nyit} (Guangyun, in Baxter's transcription), which glosses the first in this falling tone 'a holiday from work'.
\end{cousinweb}

\begin{grammarlens}[title={What you've built}]
One more word for telling what happened.
\end{grammarlens}

\subsection*{Your turn}
\begin{itemize}
\raggedright
  \item \emph{Say it:} \emph{jiàrì}
  \item \emph{Say it:} \emph{jiàrì}, once more
  \item \emph{Say it:} \emph{zhù nǐ jiàrì kuàilè} — Have a happy holiday!
  \item \emph{Recall:} say the Mandarin for to pay a visit to someone, then the Mandarin for to be reunited as a family, then say \emph{jiàrì} again
\end{itemize}

\begin{tcolorbox}[breakable,colback=teal!4,colframe=teal!35!black,title={Before you move on}]
What does \zh{假日} mean? (\textquotedblleft{}Holiday; day off\textquotedblright{}.) Say it once more.
\end{tcolorbox}
\section[fēngsú]{fēngsú — custom; social custom}

\label{lesson:ZH-C338-fengsu}



\begin{quote}
Before the new one: say the Mandarin for to be reunited as a family, then the Mandarin for a day off work.
\end{quote}

\subsection*{You'll want to know: fēngsú}
\emph{fēngsú} — \textquotedblleft{}custom; social custom\textquotedblright{}.

\begin{quote}
\emph{měi ge dìfang de fēngsú dōu bù yíyàng} — Every place has its own customs.
\end{quote}

Its characters come later; for now it is a word for the ear.

\begin{cousinweb}
\emph{fēngsú} is built from \zh{风} 'wind' and \emph{sú} 'usage', both inherited from Middle Chinese, \emph{pjuwng} and \emph{zjowk} (Guangyun, in Baxter's transcription), which glosses the second with this very word.
\end{cousinweb}

\begin{grammarlens}[title={What you've built}]
One more word for telling what happened.
\end{grammarlens}

\subsection*{Your turn}
\begin{itemize}
\raggedright
  \item \emph{Say it:} \emph{fēngsú}
  \item \emph{Say it:} \emph{fēngsú}, once more
  \item \emph{Say it:} \emph{měi ge dìfang de fēngsú dōu bù yíyàng} — Every place has its own customs.
  \item \emph{Recall:} say the Mandarin for to be reunited as a family, then the Mandarin for a day off work, then say \emph{fēngsú} again
\end{itemize}

\begin{tcolorbox}[breakable,colback=teal!4,colframe=teal!35!black,title={Before you move on}]
What does fēngsú mean? (\textquotedblleft{}Custom; social custom\textquotedblright{}.) Say it once more.
\end{tcolorbox}
